\section{Validation}
\label{sec:valid}

\begin{guidance}[Validation and automated test harness (15 marks)]
Show \textbf{passing tests for every required case} and \textbf{screenshots or a
short video} of the working system. The brief's required tests (p.~4):
CONNECT returns OK and NEW returns an empty board; a legal move changes the
board exactly as expected; an occupied-cell move returns NACK:OCCUPIED and
leaves the board unchanged; an out-of-range move returns NACK:RANGE; a move
after game over returns NACK:NOGAME (or your documented equivalent); scripted
win and draw sequences return the correct result; and, if you built an AI
opponent, it never makes an illegal move. Map each one to a named test.
\textbf{Report results on the real board}, and say clearly which results come
from a simulator, if you used one.
\end{guidance}

\subsection{\ph{Heading about your test setup}}
\ph{How the harness drives the board without human interaction (library,
e.g.~pytest~\cite{pytest} and pySerial~\cite{pyserial};
how it connects, how each test starts from a known state), and how to run it.}

\subsection{\ph{Heading that states your headline test result}}
\ph{Introduce Table~\ref{tab:required} and any further tests (robustness,
every win line, and so on). Interpret the results; do not just list them.}

\begin{table}[t]
\centering
\caption{\ph{The brief's required tests and their results on the board.}}
\label{tab:required}
\footnotesize
\begin{tabular}{@{}L{5.1cm}L{1.6cm}c@{}}
\toprule
Requirement (brief, p.~4) & Test name & Result \\
\midrule
CONNECT returns OK; NEW returns an empty board & \ph{name} & \ph{pass} \\
A legal move changes the board exactly as expected & \ph{name} & \ph{pass} \\
Occupied cell returns NACK:OCCUPIED, board unchanged & \ph{name} & \ph{pass} \\
Out-of-range move returns NACK:RANGE & \ph{name} & \ph{pass} \\
Move after game over returns NACK:NOGAME & \ph{name} & \ph{pass} \\
Scripted win and draw give the correct result & \ph{name} & \ph{pass} \\
AI never makes an illegal move (if implemented) & \ph{name} & \ph{n/a} \\
\bottomrule
\end{tabular}
\end{table}

\ph{Paste or summarise the test run output here (e.g. the final pytest summary
line), and say whether it ran against the board.}

\begin{figure}[t]
\centering
\figph{5cm}{Screenshot of your GUI connected to the board, e.g. mid-game or
just after a rejected move, with the message the user sees.}
\caption{\ph{What the screenshot shows and what to notice in it.}}
\label{fig:screenshot}
\end{figure}

\subsection{\ph{Heading about what testing revealed}}
\ph{Bugs the tests or the hardware exposed, and how you fixed them.}
